\section{Discussion}
\label{sec:discussion}

\subsection{Implications for Battery Manufacturing}
\label{sec:implications}

Our results challenge the ``purer is better'' paradigm in lithium battery manufacturing. The finding that Mg at ${\sim}1{,}000$~\ppm improves both CE and capacity retention has direct practical implications. \citet{choe2024reevaluation} estimated that relaxing purity requirements for Mg in Li$_2$CO$_3$ precursors could reduce capital expenditure by 19.4\% and CO$_2$ emissions by 9.0\% for brine-based lithium extraction. Our computational results provide mechanistic support for this approach: Mg promotes compact, low-resistivity SEI that passivates the electrode more effectively than the SEI formed in ultra-pure systems.

The threshold behavior of water impurities (neutral below ${\sim}260$~\ppm, detrimental above 500~\ppm) suggests that current electrolyte drying protocols may be overly conservative. Baakes et al.'s experimental finding that 168--260~\ppm water shows no safety difference from the reference case is consistent with our simulation results. Modest relaxation of water specifications could reduce electrolyte manufacturing costs without compromising performance.

More broadly, our finding that SEI resistivity is the dominant predictor ($r=-0.891$) provides a simple screening criterion for beneficial vs.\ detrimental impurities. Rather than treating all impurities as contaminants to be minimized, manufacturers could adopt a ``selective purity'' strategy: rigorously eliminate Fe and Cu (which increase SEI resistivity) while tolerating or deliberately introducing Mg and oxygen-containing species (which decrease it).

\subsection{Mechanistic Interpretation}
\label{sec:mechanism}

The beneficial effect of Mg can be understood through its influence on SEI structure. Mg$^{2+}$ ions promote the formation of MgO and Li$_2$O phases within the SEI, which are crystalline and ionically conductive~\citep{choe2024reevaluation}. This contrasts with the amorphous organic phases (LEDC, ROLi) that dominate in ultra-pure systems. The resulting SEI is thinner, denser, and more uniformly passivating, reducing the rate of continued electrolyte decomposition.

The dual role of water is mechanistically intuitive. At low concentrations, water promotes inorganic SEI formation through reactions that generate LiOH and subsequently Li$_2$O. At high concentrations, excess HF generation from LiPF$_6$ hydrolysis attacks the SEI and dissolves beneficial phases, while gas evolution (H$_2$, CO$_2$) disrupts the SEI mechanical integrity~\citep{baakes2023impact}.

Fe's detrimental effect arises from its redox activity. Fe$^{2+/3+}$ ions shuttle between electrodes, catalyzing electrolyte decomposition at both surfaces. At the anode, metallic Fe particles create local galvanic cells that accelerate lithium corrosion~\citep{fink2021metallic}. The resulting SEI is thick, porous, and highly resistive.

\subsection{Comparison to Literature}
\label{sec:comparison}

Our predictions agree quantitatively with published experimental results across four independent studies (\tabref{tab:comparison}).

\begin{table}[t]
    \centering
    \caption{Comparison of our computational predictions with published experimental findings. Agreement is assessed qualitatively.}
    \resizebox{\textwidth}{!}{%
    \begin{tabular}{@{}llll@{}}
        \toprule
        Finding & Our Result & Literature & Agreement \\
        \midrule
        Mg optimal ${\sim}$1~mol\% & Peak at ${\sim}$11{,}647~\ppm & Choe 2024: ${\sim}$10{,}000~\ppm & Yes \\
        Water threshold ${\sim}$260~\ppm & Neutral below 260~\ppm & Baakes 2023: 168--260~\ppm safe & Yes \\
        Inorganic SEI most stable & 127.5$^\circ$C onset & Baakes 2023: ${\sim}$128$^\circ$C & Yes \\
        Li$_2$O $>$ LiF for CE & O$_2$ species improve CE & Hobold 2024: Li$_2$O strongest correlator & Yes \\
        Fe detrimental & $-$0.33\% cap.\ ret. & Fink 2021: Fe disrupts SEI & Yes \\
        \bottomrule
    \end{tabular}
    }
    \label{tab:comparison}
\end{table}

\subsection{Limitations}
\label{sec:limitations}

We acknowledge several limitations of this study:

\para{Simulation fidelity.}
The \pybamm SPMe model captures bulk electrochemical behavior but does not explicitly resolve impurity-specific mechanisms such as Fe dissolution-plating, Mg site substitution in cathode crystal structures, or spatially heterogeneous SEI growth. The CE differences we observe (0.001--0.003\%) are small relative to the model's numerical precision, as reflected in the non-significant ANOVA on CE. Capacity retention is a more robust metric in our framework.

\para{Parameter calibration.}
Impurity effects were modeled as multiplicative scaling factors (5--50\% per 100~\ppm) on base SEI parameters. While the qualitative directions are well-supported by literature, the exact scaling relationships require experimental calibration. Physical validation with controlled impurity additions is essential before translating our predictions to manufacturing guidelines.

\para{Single cell chemistry.}
All simulations used Chen 2020 parameters for NMC/graphite chemistry. Impurity effects may differ for LFP, NCA, or Li metal anode systems due to different SEI formation mechanisms and operating potentials.

\para{Short cycling window.}
Our 100-cycle simulations capture early degradation trends but may miss long-term effects. Beneficial impurities that slightly accelerate SEI growth could become detrimental over 500+ cycles if the SEI grows too thick. Extended simulations are needed to assess long-term trajectories.

\para{No stochastic variation.}
Each simulation is deterministic. Real cells exhibit cell-to-cell variability (typically 0.1--1\% in CE) that could mask or amplify the small impurity effects we observe. Multi-cell experimental validation with statistical replication is critical.

\para{Cross-study variability.}
Literature data points were compiled from studies with different cell formats, protocols, and CE definitions, introducing systematic uncertainty into our dose-response fits.
